\begin{table*}[h!]
    \centering
\resizebox{\textwidth}{!}{%
\begin{tabular}{llcllll}
\hline
\multirow{2}{*}{\textbf{Direction}} &
  \multirow{2}{*}{\textbf{Metric}} &
  \multirow{2}{*}{\textbf{\begin{tabular}[c]{@{}c@{}}Reference\\ Summary\end{tabular}}} &
  \multicolumn{1}{c}{\multirow{2}{*}{\textbf{\begin{tabular}[c]{@{}c@{}}Video\\ Used\end{tabular}}}} &
  \multicolumn{3}{c}{\textbf{MLVU(VS)-Eval}} \\ \cline{5-7} 
 &
   &
   &
  \multicolumn{1}{c}{} &
  \multicolumn{1}{c}{\textbf{Kendall's $\tau_b$}} &
  \multicolumn{1}{c}{\textbf{Kendall's $\tau_c$}} &
  \multicolumn{1}{c}{\textbf{Spearman's $\rho$}} \\ \hline
\multirow{5}{*}{Rule-based}       & BLEU-1                   & \checkmark            &                               & 0.0217 (0.7803) & 0.0219 (0.7803) & 0.0431 (0.7045) \\
                                  & BLEU-2                   & \checkmark            &                               & 0.0673 (0.3865) & 0.0680 (0.3865) & 0.1045 (0.3561) \\
                                  & BLEU-3                   & \checkmark            &                               & 0.1994 (0.0152) & 0.1836 (0.0152) & 0.2836 (0.0108) \\
                                  & BLEU-4                   & \checkmark            &                               & 0.2113 (0.0155) & 0.1578 (0.0155) & 0.2798 (0.0119) \\
                                  & ROUGE-L                  & \checkmark            &                               & 0.1094 (0.1593) & 0.1104 (0.1593) & 0.1571 (0.1639) \\ \hline
\multirow{5}{*}{Similarity-based} & METEOR                   & \checkmark            &                               & 0.2663 (0.0006) & 0.2689 (0.0006) & 0.3822 (0.0005) \\
                                  & CIDEr                    & \checkmark            &                               & 0.2401 (0.0015) & 0.2435 (0.0014) & 0.3612 (0.0012) \\
                                  & SPICE                    & \checkmark            &                               & 0.2287 (0.0021) & 0.2312 (0.0020) & 0.3489 (0.0019) \\
                                  & BERTscore                & \checkmark            &                               & 0.1415 (0.0688) & 0.1428 (0.0688) & 0.1987 (0.0773) \\
                                  & Video-Summary Similarity & \multicolumn{1}{l}{} & \multicolumn{1}{c}{\checkmark}                     & 0.0278 (0.0517) & 0.0266 (0.0522) & 0.0401 (0.0549) \\ \hline
\multirow{2}{*}{LLM-based}        & MLVU                     & \checkmark            &                               & 0.5284 (0.0000) & 0.5309 (0.0000) & 0.6738 (0.0000) \\
                                  & Video-LMM Judge          & \multicolumn{1}{l}{} & \multicolumn{1}{c}{\checkmark} & 0.5376 (0.0000) & 0.5441 (0.0000) & 0.6810 (0.0000) \\ \hline
QA-based                          & QEVA(Ours)               & \multicolumn{1}{l}{} & \multicolumn{1}{c}{\checkmark} & \textbf{0.6465 (0.0000)} & \textbf{0.6407 (0.0000)} & \textbf{0.7326 (0.0000)} \\ \hline
\end{tabular}%
}
\caption{Comparative evaluation results of QEVA and existing summarization evaluation metrics on the MLVU(VS)-Eval benchmark. We report correlations with human judgments using Kendall's $\tau_b$, $\tau_c$, and Spearman's $\rho$. The metrics are categorized into several groups: rule-based n-gram metrics (BLEU variants, ROUGE-L), similarity-based metrics (METEOR, CIDEr, SPICE, BERTScore, Video-Summary Similarity), LLM-based metrics (MLVU, Video-LMM Judge), and our proposed QA-based metric (QEVA). Reference usage and video modality usage for each metric are also indicated.
\textbf{Takeaway:} QEVA consistently achieves significantly higher correlation with human judgments compared to existing metrics, demonstrating its effectiveness as a reference-free and multimodal evaluation metric.}
    \label{tab:correlation-results}
    \end{table*}
% Please add the following required packages to your document preamble:
% \usepackage{multirow}
% \usepackage{graphicx}
\begin{table*}[h!]
\centering
\resizebox{\textwidth}{!}{%
\begin{tabular}{l|cccc|cccc}
\hline
\multirow{2}{*}{\textbf{Metric}} &
  \multicolumn{4}{c|}{\textbf{Kendalls' $\tau_b$}} &
  \multicolumn{4}{c}{\textbf{Kendalls' $\tau_c$}} \\ \cline{2-9} 
 &
  \textbf{Qwen2.5-VL-7B} &
  \textbf{InternVL3-8B} &
  \textbf{Video-LLaVA-7B} &
  \textbf{GPT-4o} &
  \textbf{Qwen2.5-VL-7B} &
  \textbf{InternVL3-8B} &
  \textbf{Video-LLaVA-7B} &
  \textbf{GPT-4o} \\ \hline
BLEU-1                   & 0.1662 & 0.0870  & 0.0758          & 0.0053 & 0.1669  & 0.0867  & 0.0758          & 0.0054 \\
BLEU-2                   & 0.2735 & 0.1247  & 0.1457          & 0.1333 & 0.2746  & 0.1246  & 0.1458          & 0.1339 \\
BLEU-3                   & 0.2914 & 0.1380  & 0.1370          & 0.3293 & 0.2800  & 0.1300  & 0.1108          & 0.3088 \\
BLEU-4                   & 0.1654 & 0.1130  & \textbf{0.3565} & 0.1717 & 0.1444  & 0.0914  & 0.2025          & 0.1463 \\
ROUGE-L                  & 0.4022 & 0.2332  & 0.1573          & 0.3423 & 0.4038  & 0.2329  & 0.1575          & 0.3429 \\ \hline
METEOR                   & 0.2628 & 0.1573  & 0.3555          & 0.4107 & 0.2638  & 0.1571  & \textbf{0.3558} & 0.4125 \\
BERTscore                & 0.5202 & 0.0163  & 0.0701          & 0.3147 & 0.5223  & 0.0163  & 0.0700          & 0.3161 \\
Video-Summary Similarity & 0.0012 & -0.0301 & 0.2277          & 0.2417 & -0.0218 & -0.0119 & 0.2117          & 0.1472 \\ \hline
MLVU                     & 0.3950 & 0.3260  & -0.1254         & 0.2361 & 0.3900  & 0.3208  & -0.1067         & 0.2392 \\
Video-LMM Judge          & 0.4138 & 0.3358  & 0.2056          & 0.2178 & 0.2005  & 0.3123  & 0.1233          & 0.2198 \\ \hline
\textbf{QEVA(Ours)} &
  \textbf{0.4509} &
  \textbf{0.4268} &
  0.1450 &
  \textbf{0.4262} &
  \textbf{0.4500} &
  \textbf{0.4114} &
  0.1000 &
  \textbf{0.4222} \\ \hline
\end{tabular}%
}
\caption{Model-wise correlation analysis of summarization evaluation metrics on the MLVU(VS)-Eval dataset. Correlations (measured by Kendall's $\tau_b$ and $\tau_c$) are reported separately for summaries generated by four representative Video-LMM models: Qwen2.5-VL-7B, InternVL3-8B, Video-LLaVA-7B, and GPT-4o.
\textbf{Takeaway:} QEVA exhibits stable and positive correlations with human judgments across diverse Video-LMM summarization models, highlighting its robustness and generalizability.}
\label{tab:model-wise correlation}
\end{table*}

\section{Experiments}
\label{sec:experiments}

We conduct comprehensive experiments to validate the effectiveness of our proposed QEVA metric. We first introduce \textbf{MLVU(VS)-Eval}, a novel benchmark dataset for evaluating video-to-text summarization metrics (\S\ref{sec:dataset}). We then compare QEVA against existing metrics in terms of correlation with human judgments (\S\ref{sec:correlation}). Additionally, we analyze model-wise performance of QEVA (\S\ref{sec:modelwise-corr}), and demonstrate the robustness of QEVA via ablation studies (\S\ref{sec:ablation}).

\subsection{MLVU(VS)-Eval: A Novel Benchmark for Evaluating Video Summarization Metrics} \label{sec:dataset}

Existing evaluation datasets for video-to-text summarization lack human annotations, limiting accurate metric evaluation. To address this, we propose \textbf{MLVU(VS)-Eval}, a novel human-annotated dataset built upon the MLVU benchmark~\citep{zhou2024mlvu}.

We select 200 video clips (average length $\sim$15 minutes) from the MLVU Video Summarization task. Each video has a human-written reference summary. We generate candidate summaries using four widely-used Video-LMM models: GPT-4o, InternVL3-8B, Qwen2.5-VL-7B, and Video-LLaVA-7B, resulting in 800 candidate summaries. For our evaluation, we recruited 20 annotators (comprising graduate and undergraduate students) to assess each summary based on three criteria: \textit{Coverage}, \textit{Factuality}, and \textit{Chronology}, using a 5-point Likert scale. Each summary received evaluations from two independent annotators. The inter-annotator agreement, measured using Krippendorff's $\alpha$, is 0.68, indicating substantial reliability of the annotations.

\subsection{Correlation with Human Judgments} \label{sec:correlation}

We compare QEVA with various existing evaluation metrics in terms of correlation with human judgments.

\paragraph{Baseline Metrics.} 
We compare QEVA with several representative categories of evaluation metrics: reference-based n-gram overlap metrics, embedding-based similarity metrics, and LLM-based evaluation approaches. For the n-gram and embedding-based metrics, we report results using widely adopted methods in each category. Additionally, we introduce two novel baselines: (1) employing a Video-LMM as a direct judge by prompting it with our human annotation guidelines, and (2) computing multimodal embedding similarity between the original video and generated summary using a state-of-the-art multimodal encoder\footnote{We use LanguageBind~\citep{zhu2023languagebind} for multimodal embedding in this paper.}.

\paragraph{Results and Analysis.} Table~\ref{tab:correlation-results} presents correlations measured by Kendall's $\tau_b$, $\tau_c$, and Spearman's $\rho$. QEVA consistently achieves significantly higher correlations than all existing metrics, highlighting its superior alignment with human evaluations, especially notable given QEVA is a reference-free metric.

\subsection{Model-wise Correlation Analysis} \label{sec:modelwise-corr}

To further examine QEVA's consistency, we analyze correlation results separately for each Video-LMM summarization model. Table~\ref{tab:model-wise correlation} shows that QEVA consistently yields positive and stable correlations across different summarization models. In contrast, some baseline metrics even show negative or inconsistent correlations for certain models.

We observe relatively lower correlations for Video-LLaVA-generated summaries. Upon qualitative analysis, we find these summaries are generally shorter and of lower quality, making it challenging for QEVA to generate meaningful questions and answers.

\subsection{Ablation Studies} \label{sec:ablation}

We perform ablation studies to analyze QEVA's robustness and generalizability.

\paragraph{Evaluation Criteria-wise Ablation.} 
We separately calculate correlations between human judgments and QEVA scores for each evaluation criterion (Coverage, Factuality, Chronology). As shown in Table~\ref{tab:criterion-ablation}, QEVA demonstrates strong correlations across all individual criteria, validating its capability to accurately capture different aspects of summarization quality.

\begin{table}[h!]
\centering
\begin{tabularx}{\columnwidth}{l|>{\centering\arraybackslash}X >{\centering\arraybackslash}X >{\centering\arraybackslash}X}
\hline
\textbf{Criteria} & \textbf{$\tau_b$} & \textbf{$\tau_c$} & \textbf{$\rho$} \\ 
\hline
\textbf{Coverage}   & 0.6005 & 0.5872 & 0.7349 \\
\textbf{Factuality} & 0.5731 & 0.5586 & 0.7123 \\
\textbf{Chronology} & 0.5610 & 0.5479 & 0.6984 \\
\hline
\end{tabularx}
\caption{Ablation study of QEVA across the three individual evaluation criteria (Coverage, Factuality, Chronology). Correlations with human judgments (measured by Kendall's $\tau_b$, $\tau_c$, and Spearman's $\rho$) are reported separately for each dimension. 
\textbf{Takeaway:} QEVA demonstrates strong and consistent correlations across all individual evaluation criteria, validating its comprehensive and multidimensional evaluation approach.}
\label{tab:criterion-ablation}
\end{table}
% Please add the following required packages to your document preamble:
% \usepackage{graphicx}
% \usepackage[normalem]{ulem}
% \useunder{\uline}{\ul}{}
\begin{table}[t]
\centering
\resizebox{\columnwidth}{!}{%
\begin{tabular}{l|ccc}
\hline
\textbf{Video-LMM + LLM}         & \textbf{$\tau_b$} & \textbf{$\tau_c$} & \textbf{$\rho$}  \\ \hline
Qwen2.5-VL + LLaMA-3.1           & 0.6211            & 0.6152            & 0.7012          \\
Qwen2.5-VL + Gemma-3             & 0.6075            & 0.6021            & 0.7120          \\
InternVL3 + LLaMA-3.1            & 0.5988            & 0.5933            & 0.6761          \\
InternVL3 + Gemma-3              & 0.5827            & 0.5780            & 0.6894          \\
\textbf{Gemini-1.5-Pro + GPT-4o} & \textbf{0.6465}   & \textbf{0.6407}   & \textbf{0.7326} \\ \hline
\end{tabular}%
}
\caption{Ablation study examining the robustness of QEVA to different combinations of Video-LMM and LLM models compared to the default setting (Gemini-1.5 Pro + GPT-4o). We report correlations with human annotations (Kendall's $\tau_b$, $\tau_c$, and Spearman's $\rho$) when replacing original models with alternative open-source models (Qwen2.5-VL, InternVL3, LLaMA-3.1, Gemma-3).
\textbf{Takeaway:} QEVA maintains high correlations with human judgments even when using alternative open-source models, indicating practical applicability and cost-effectiveness without relying solely on costly API-based models.}
\label{tab:component-ablation}
\end{table}

\paragraph{Robustness to Different Video-LMM and LLM Models.} 
We further examine QEVA's internal robustness by replacing the original Video-LMM (Gemini-1.5 Pro) and LLM (GPT-4o) components with alternative models. Specifically, we test open-source Video-LMM variants (QwenVL, InternVL) and LLM variants (LLaMA, Gemma). Table~\ref{tab:component-ablation} indicates that QEVA maintains high correlations even with open-source alternatives. This demonstrates QEVA's practical applicability and cost-effectiveness by alleviating reliance on costly API-based models.

